\part{Condicionales}


\section{Booleanos}

Otro tipo de dato es el booleano, que representa valores de verdad:
\begin{table}[h]
\centering
\footnotesize
\begin{tabular}{|l|l|l|p{3.5cm}|l|}
\hline
\textbf{Tipo de dato} & \textbf{Nombre (inglés)} & \textbf{Traducción} & \textbf{Descripción} & \textbf{E.g.} \\
\hline
\mintinline{python}|bool| & \textit{boolean} & booleano & Representa verdadero o falso & \mintinline{python}|True|, \mintinline{python}|False| \\
\hline
\end{tabular}
\normalsize
\caption{Tipo de dato booleano.}
\label{tab:tipo-dato-booleano}
\end{table}

El término \emph{booleano} alude al álgebra de Boole, un sistema que pretendía modelar el razonamiento lógico mediante herramientas algebraicas. Resultó en una estructura algebraica que utiliza solo dos valores. En Python, los dos valores son \texttt{True} y \texttt{False} (nótese que se escriben con mayúscula inicial), pero conceptualmente podrían representarse en cualquier forma, como \texttt{1} y \texttt{0}.

\subsection{Operadores relacionales}

Una forma de generar valores booleanos a partir de valores que no son booleanos es mediante los operadores relacionales, que comparan dos valores y devuelven un valor booleano:
\begin{table}[h]
\centering
\begin{tabular}{|l|c|c|c|c|}
\hline
\textbf{Operador} & \textbf{Símbolo} & \textbf{Ejemplo matemático} & \textbf{Ejemplo en Python} & \textbf{Resultado} \\
\hline
Igual a & \mintinline{python}|==| & $1 = 2$ & \mintinline{python}|1 == 2| & \mintinline{python}|False| \\
\hline
Distinto de & \mintinline{python}|!=| & $5 \neq 3$ & \mintinline{python}|5 != 3| & \mintinline{python}|True| \\
\hline
Mayor que & \mintinline{python}|>| & $3 > 4$ & \mintinline{python}|3 > 4| & \mintinline{python}|False| \
\hline
Menor que & \mintinline{python}|<| & $2 < 6$ & \mintinline{python}|2 < 6| & \mintinline{python}|True| \\
\hline
Mayor o igual que & \mintinline{python}|>=| & $5 \geq 5$ & \mintinline{python}|5 >= 5| & \mintinline{python}|True| \\
\hline
Menor o igual que & \mintinline{python}|<=| & $5 \leq 4$ & \mintinline{python}|5 <= 4| & \mintinline{python}|False| \\
\hline
\end{tabular}
\caption{Operadores relacionales en Python.}
\label{tab:operadores-relacionales}
\end{table}

\begin{tip}
    Note la diferencia entre el operador de asignación \texttt{=} y el operador de igualdad \texttt{==}. El primero se usa para asignar un valor a una variable, mientras que el segundo se usa para comparar dos valores.
\end{tip}

\subsubsection{Operadores de identidad}

Otra forma de generar valores booleanos es mediante los operadores de identidad, que comparan la identidad de dos objetos en memoria. En Python, los operadores de identidad son \texttt{is} e \texttt{is not}:

\begin{table}[h]
\centering
\begin{tabular}{|l|c|c|c|}
\hline
\textbf{Operador} & \textbf{Símbolo} & \textbf{Ejemplo en Python} & \textbf{Resultado} \\
\hline
Es idéntico a & \mintinline{python}|is| & \mintinline{python}|type(.1) is float| & \mintinline{python}|True| \\
\hline
No es idéntico a & \mintinline{python}|is not| & \mintinline{python}|5 is not None| & \mintinline{python}|True| \\
\hline
\end{tabular}
\caption{Operadores de identidad en Python.}
\label{tab:operadores-identidad}
\end{table}

En Python, dos cosas son la misma en memoria si son el mismo objeto; sin embargo, la programación orientada a objetos escapa el alcance de estos apuntes.

Lo que hay que saber es que se debe usar \mintinline{python}|is| e \mintinline{python}|is not| para comparar con \mintinline{python}|None|, al igual que para comparar con tipos de dato cuando se usa la función \texttt{type} (e.g., \mintinline{python}|type(x) is int|). Para comparar otros valores, se suele usar \texttt{==} y \texttt{!=}.

\subsection{Operadores lógicos}

Los operadores lógicos generan valores booleanos a partir de otros valores booleanos:

\begin{table}[h]
\centering
\begin{tabular}{|l|c|c|c|c|}
\hline
\textbf{Operador} & \textbf{Símbolo} & \textbf{Ejemplo en español} & \textbf{Ejemplo matemático} & \textbf{Ejemplo en Python} & \textbf{Resultado} \\
\hline
Negación & \mintinline{python}|not| & No $p$ & $\neg (p)$ & \mintinline{python}|not p| & \mintinline{python}|False| \\
\hline
Conjunción & \mintinline{python}|and| & $p$ y $q$ & $p \land q$ & \mintinline{python}|p and q| & \mintinline{python}|True| si y solo si $p$ y $q$ son \texttt{True}\\
\hline
Disyunción & \mintinline{python}|or| & $p$ o $q$ & $p \lor q$ & \mintinline{python}|p or q| & \mintinline{python}|True| \\
\hline
\end{tabular}
\caption{Operadores lógicos en Python.}
\label{tab:operadores-logicos}
\end{table}

Uso  \mintinline{python}|p| y \mintinline{python}|q| como los nombres de las variables porque es convencional en lógica usar letras minúsculas para representar proposiciones, que son afirmaciones que pueden ser verdaderas o falsas. En este caso, \mintinline{python}|p| y \mintinline{python}|q| son valores de tipo \texttt{bool}.


\section{Booleanos}

\subsection{Operadores relacionales}

% Los operadores de datos relacionales y cómo generan valores bool a partir de valores que no son booleanos. (L2).

\subsubsection{Operadores de identidad}

% Los operadores de identidad que, como mencioné, no son un tema central del curso. (L2).

\subsection{Operadores lógicos}

% Los operadores lógicos, cómo funcionan y cómo interactúan entre sí. (L2).

\subsubsection{Tablas de verdad}

% Las tablas de verdad para los tres operadores. (L3).

\section{Estructura de control condicional}

La estructura de control condicional en Python utiliza la palabra \texttt{if}, que traduce al español como \textit{si}. Su sintaxis es:
\begin{sintaxis}
if \placeholder{valor de tipo \texttt{bool}}:\\
\phantom{xxxx}\placeholder{bloque de código}
\end{sintaxis}
Por ejemplo:
\begin{minted}{python}
if True:
    print("Hello, World!")
\end{minted}

El bloque de código se ejecuta si la condición es verdadera.



% La estructura de control condicional en Python.
% Cómo se escribe.
% Qué hace if y cuándo se ejecuta el bloque de código en su cuerpo.
% Qué hace elif, cuándo se ejecuta el bloque de código en su cuerpo y cuándo tiene sentido escribir uno (o varios) elifs.
% Qué hace else, cuándo se ejecuta el bloque de código en su cuerpo y cuándo tiene sentido usarlo.

\subsection{Cascadas de condicionales vs condicionales en secuencia}
% La diferencia entre una cascada de condicionales, donde evaluamos condiciones mutuamente excluyentes usando una sola estructura de control condicional, y condicionales en secuencia, donde tenemos varias estructuras de control condicionales independientes, una después de la otra.

\subsection{Expresión ternaria}

\section{Abordar problemas de programación con condicionales}

% Hoy vimos herramientas conceptuales para abordar problemas de programación que requieren lógica:

% Diagramas de flujo, que son (a este nivel) una representación gráfica de una secuencia de decisiones. Vimos que se pueden traducir a código directamente.
% Simplificación lógica, que consiste en representar el enunciado del problema con una o más expresiones con variables y operadores lógicos. Eso permite simplificar las condiciones usando las propiedades de los operadores lógicos.
% Vimos ejemplos de problemas complejos que se vuelven más sencillos usando esas estrategias. De esas herramientas espero que se queden con dos cosas:

% Frecuentemente es mejor pensar el problema y hacer una representación de la solución sin usar código, en papel, antes de empezar a escribir código. Esto es un consejo no solo para esta materia, sino para todos los niveles de programación (y de resolución de problemas).
% A menudo resulta útil razonar los problemas de forma “top-down”, es decir, pensando primero en solucionar el problema más grande y luego en solucionar los problemas chiquitos que van apareciendo.

\section{Diccionarios}

